%----------------------------------------------------------------------------------------
% GONIOMETRISCHE VERGELIJKINGEN. Alle opgaven zijn eigen werk.
%
% Opbouw: Start, leerstof (theorie, werkwijze of uitleg, voorbeeld, korte oefening), en op een
% nieuwe pagina de Opdrachten (door elkaar, stijgende moeilijkheid), Koppel en Test jezelf.
% Labels van dit hoofdstuk beginnen met h4-.
%----------------------------------------------------------------------------------------

\mychapter{4}{Goniometrische vergelijkingen}

\lkr{Leerplan Basis: LPD 11-13 en K2 (goniometrische functies), het deel over de cyclometrische functies en goniometrische vergelijkingen. De algemene oplossing met $k \in \Z$ is voor veel leerlingen abstract: laat ze telkens ook de oplossingen in een interval opschrijven, en controleren op de eenheidscirkel. Het rekentoestel geeft maar één oplossing; de tweede vinden ze met de verwante hoeken.}

Wanneer is het water in de haven diep genoeg? Wanneer is de gondel hoger dan 100 m? Zulke vragen leiden tot een vergelijking waarin de onbekende in een sinus of cosinus staat. Omdat de functies periodiek zijn, hebben zulke vergelijkingen meestal oneindig veel oplossingen.

We gaan in drie stappen te werk. Eerst zoeken we bij een gegeven sinus, cosinus of tangens één hoek: dat doen de cyclometrische functies. Dan vinden we met de eenheidscirkel alle oplossingen. Ten slotte herleiden we moeilijkere vergelijkingen tot zulke eenvoudige.

\begin{oefeningen}[Start]
\oef{h4-start-1}

\oef{h4-verken}
\end{oefeningen}

%========================================================================================
\section{Cyclometrische functies}

De sinus is niet injectief: $\sin\frac{\pi}{6} = \sin\frac{5\pi}{6}$, en ook $\sin\frac{13\pi}{6}$ is $\frac12$. Een functie die niet injectief is, heeft geen inverse. Beperk je het domein tot een stuk waarop de functie elke waarde maar één keer aanneemt, dan heeft ze daar wel een inverse. Haar grafiek is het spiegelbeeld om de rechte $y = x$. Zo krijgen we drie nieuwe functies: de cyclometrische functies.

\subsection{Boogsinus}

Op $\bigl[-\frac{\pi}{2}, \frac{\pi}{2}\bigr]$ stijgt de sinus van $-1$ tot $1$, en neemt ze elke waarde in $[-1, 1]$ precies één keer aan.

\begin{theorie}
De \textbf{\textcolor{accent}{boogsinus}} is de inverse van de sinus op $\bigl[-\frac{\pi}{2}, \frac{\pi}{2}\bigr]$. Ze heeft domein $[-1, 1]$ en bereik $\bigl[-\frac{\pi}{2}, \frac{\pi}{2}\bigr]$:

\vspace{1mm}
\qquad $\formule{\bgsin a = \theta \iff \sin\theta = a \text{ en } -\tfrac{\pi}{2} \leqslant \theta \leqslant \tfrac{\pi}{2}}$
\end{theorie}

\naast{%
De naam komt van \textit{boogsinus}: $\bgsin a$ is de boog (de hoek in radialen) waarvan de sinus $a$ is. Op het rekentoestel heet de toets $\sin^{-1}$; zet het toestel in radialen.

De grafiek van $\bgsin$ is het spiegelbeeld van het stijgende stuk van de sinus om de rechte $y = x$.

Let op: $\bgsin a$ is één hoek, uit het bereik $\bigl[-\frac{\pi}{2}, \frac{\pi}{2}\bigr]$. Zo is $\sin\frac{5\pi}{6} = \frac12$, maar $\bgsin\frac12 = \frac{\pi}{6}$, want $\frac{5\pi}{6}$ ligt niet in het bereik.}{%
\begin{tikzpicture}[x=.85cm,y=.85cm,font=\scriptsize]
  \draw[->,grijs] (-1.9,0) -- (2,0) node[right,text=black] {$x$};
  \draw[->,grijs] (0,-1.9) -- (0,2) node[above,text=black] {$y$};
  \draw[grijs,dashed] (-1.8,-1.8) -- (1.8,1.8);
  \draw[accent,thick,dashed] plot[domain=-1.5708:1.5708,samples=50,smooth] (\x,{sin(\x r)});
  \draw[accent,thick] plot[domain=-1:1,samples=80,smooth] (\x,{asin(\x)*pi/180});
  \node[accent,right] at (1,1.5708) {$\bgsin x$};
  \node[accent,below right] at (1.4,.9) {$\sin x$};
  \node[left] at (0,1.5708) {$\frac{\pi}{2}$};
  \node[left] at (0,-1.5708) {$-\frac{\pi}{2}$};
\end{tikzpicture}}

\voorbeeld{}{Bepaal $\bgsin 0{,}8$ op vier decimalen, en geef nog een hoek in $[0, 2\pi[$ met sinus $0{,}8$.\label{vb:h4-bgsin}}{%
\uitleg{$0{,}8$ is geen waarde uit de tabel, dus we gebruiken het rekentoestel. Het geeft de hoek uit het bereik van $\bgsin$. Een tweede hoek met dezelfde sinus is de supplementaire hoek.}
\begin{stappenbreed}
\bgsin 0{,}8 \approx 0{,}9273 & Rekentoestel in radialen: $\sin^{-1}(0{,}8)$. \\[1ex]
\sin 0{,}9273 \approx 0{,}8 \text{ en } -\tfrac{\pi}{2} \leqslant 0{,}9273 \leqslant \tfrac{\pi}{2} & Controle met de definitie. \\[1ex]
\pi - 0{,}9273 \approx 2{,}2143 & Supplementaire hoeken hebben dezelfde sinus. Deze hoek ligt niet in het bereik van $\bgsin$.
\end{stappenbreed}}

\begin{oefeningen}
\oef{h4-bgsin-k}
\end{oefeningen}

\subsection{Boogcosinus}

De cosinus is niet injectief op $\bigl[-\frac{\pi}{2}, \frac{\pi}{2}\bigr]$, want daar is $\cos(-\theta) = \cos\theta$. Op $[0, \pi]$ daalt ze wel van $1$ tot $-1$, en neemt ze elke waarde in $[-1, 1]$ precies één keer aan.

\begin{theorie}
De \textbf{\textcolor{accent}{boogcosinus}} is de inverse van de cosinus op $[0, \pi]$. Ze heeft domein $[-1, 1]$ en bereik $[0, \pi]$:

\vspace{1mm}
\qquad $\formule{\bgcos a = \theta \iff \cos\theta = a \text{ en } 0 \leqslant \theta \leqslant \pi}$
\end{theorie}

\naast{%
Op de eenheidscirkel is $\bgcos a$ de hoek van het punt met $x$-coördinaat $a$ in de bovenste helft van de cirkel. Voor een negatieve $a$ ligt dat punt in het tweede kwadrant, en is $\bgcos a$ een stompe hoek. Dat is anders dan bij $\bgsin$, die voor een negatieve $a$ een negatieve hoek geeft.

De grafiek daalt van $(-1, \pi)$ naar $(1, 0)$.}{%
\begin{tikzpicture}[x=.75cm,y=.75cm,font=\scriptsize]
  \draw[->,grijs] (-1.5,0) -- (3.6,0) node[right,text=black] {$x$};
  \draw[->,grijs] (0,-1.4) -- (0,3.6) node[above,text=black] {$y$};
  \draw[grijs,dashed] (-1.3,-1.3) -- (3.4,3.4);
  \draw[accent,thick,dashed] plot[domain=0:3.1416,samples=50,smooth] (\x,{cos(\x r)});
  \draw[accent,thick] plot[domain=-1:1,samples=80,smooth] (\x,{acos(\x)*pi/180});
  \node[accent,above right] at (-1,3.1416) {$\bgcos x$};
  \node[accent,below] at (2.6,-.9) {$\cos x$};
  \node[left] at (0,3.1416) {$\pi$};
  \node[below] at (3.1416,0) {$\pi$};
\end{tikzpicture}}

\voorbeeld{}{Bepaal exact $\bgcos\bigl(-\frac{\sqrt2}{2}\bigr)$.\label{vb:h4-bg}}{%
\uitleg{Zoek de hoek in $[0, \pi]$ met cosinus $-\frac{\sqrt2}{2}$. In dat interval is de cosinus negatief in het tweede kwadrant.}
\begin{stappenbreed}
\cos\theta = -\tfrac{\sqrt2}{2},\ 0 \leqslant \theta \leqslant \pi & Definitie. \\[1ex]
\theta = \pi - \tfrac{\pi}{4} = \tfrac{3\pi}{4} & Supplementair met $\frac{\pi}{4}$: $\cos\bigl(\pi - \frac{\pi}{4}\bigr) = -\cos\frac{\pi}{4} = -\frac{\sqrt2}{2}$.
\end{stappenbreed}}

\begin{oefeningen}
\oef{h4-bgcos-k}
\end{oefeningen}

\subsection{Boogtangens}

Op $\bigl]-\frac{\pi}{2}, \frac{\pi}{2}\bigr[$, één tak van de tangens, stijgt de tangens van $-\infty$ naar $+\infty$. Ze neemt er elke reële waarde precies één keer aan.

\begin{theorie}
De \textbf{\textcolor{accent}{boogtangens}} is de inverse van de tangens op $\bigl]-\frac{\pi}{2}, \frac{\pi}{2}\bigr[$. Ze heeft domein $\R$ en bereik $\bigl]-\frac{\pi}{2}, \frac{\pi}{2}\bigr[$:

\vspace{1mm}
\qquad $\formule{\bgtan a = \theta \iff \tan\theta = a \text{ en } -\tfrac{\pi}{2} < \theta < \tfrac{\pi}{2}}$
\end{theorie}

\naast{%
Bij het spiegelen om $y = x$ worden de verticale asymptoten van de tangens horizontale asymptoten $y = \pm\frac{\pi}{2}$ van de boogtangens. Hoe groot $a$ ook is, $\bgtan a$ blijft kleiner dan $\frac{\pi}{2}$: de rechte $OP$ wordt steiler, maar nooit verticaal.}{%
\begin{tikzpicture}[x=.5cm,y=.62cm,font=\scriptsize]
  \draw[->,grijs] (-4.4,0) -- (4.6,0) node[right,text=black] {$x$};
  \draw[->,grijs] (0,-2.1) -- (0,2.3) node[above,text=black] {$y$};
  \draw[grijs,dashed] (-4.2,1.5708) -- (4.2,1.5708);
  \draw[grijs,dashed] (-4.2,-1.5708) -- (4.2,-1.5708);
  \draw[accent,thick] plot[domain=-4.2:4.2,samples=80,smooth] (\x,{atan(\x)*pi/180});
  \node[left] at (0,1.85) {$\frac{\pi}{2}$};
  \node[right] at (0,-1.85) {$-\frac{\pi}{2}$};
  \node[accent,above] at (3,1.05) {$\bgtan x$};
\end{tikzpicture}}

\voorbeeld{}{Bepaal $\bgtan 5$ op vier decimalen. Waarom ligt het resultaat zo dicht bij $\frac{\pi}{2}$?\label{vb:h4-bgtan}}{%
\uitleg{Het rekentoestel geeft de hoek uit $\bigl]-\frac{\pi}{2}, \frac{\pi}{2}\bigr[$. Daarna denken we aan de rechte $OP$ met richtingscoëfficiënt $5$.}
\begin{stappenbreed}
\bgtan 5 \approx 1{,}3734 & Rekentoestel in radialen: $\tan^{-1}(5)$. \\[1ex]
\tfrac{\pi}{2} \approx 1{,}5708 & Ter vergelijking. \\[1ex]
\text{rico } 5 & De rechte $OP$ is steil: ze stijgt $5$ per $1$ naar rechts. De hoek met de $x$-as is dus bijna een rechte hoek.
\end{stappenbreed}}

\nujij{Bepaal exact $\bgsin\bigl(-\frac12\bigr)$ en $\bgtan 1$.}{$-\frac{\pi}{6}$, want $\sin\bigl(-\frac{\pi}{6}\bigr) = -\frac12$ en $-\frac{\pi}{6}$ ligt in $\bigl[-\frac{\pi}{2}, \frac{\pi}{2}\bigr]$. $\bgtan 1 = \frac{\pi}{4}$.}

\needspace{5\baselineskip}
\begin{oefeningen}
\oef{h4-bg}
\end{oefeningen}

%========================================================================================
\section{Elementaire vergelijkingen}

Een \textbf{elementaire vergelijking} is van de vorm $\sin x = a$, $\cos x = a$ of $\tan x = a$. Een cyclometrische functie geeft één oplossing. De eenheidscirkel toont waar de andere liggen.

\subsection{De vergelijking $\sin x = a$}

\naast{%
Een horizontale rechte $y = a$ (met $-1 < a < 1$) snijdt de eenheidscirkel in twee punten, elkaars spiegelbeeld om de $y$-as. Het rechtse punt hoort bij de hoek $\theta = \bgsin a$, het linkse bij de supplementaire hoek $\pi - \theta$. Elke hoek die bij een van die punten hoort, is een oplossing: daar komen dus nog veelvouden van $2\pi$ bij.}{%
\begin{tikzpicture}[scale=1.3,font=\scriptsize]
  \eenheidscirkel
  \draw[accent] (-1.2,.45) -- (1.25,.45) node[right] {$y = a$};
  \draw[accent] (0,0) -- (26.7:1) (0,0) -- (153.3:1);
  \fill[accent] (26.7:1) circle (.035);
  \fill[accent] (153.3:1) circle (.035);
  \node[accent] at (12:.55) {$\theta$};
  \node[accent,above] at (153.3:1.05) {$\pi - \theta$};
\end{tikzpicture}}

\begin{theorie}
\begin{stelling}{De vergelijking $\sin x = a$}
Voor $-1 \leqslant a \leqslant 1$ en met $k \in \Z$:

\vspace{1mm}
\qquad $\formule{\sin x = a \iff x = \bgsin a + k \cdot 2\pi \ \text{ of } \ x = \pi - \bgsin a + k \cdot 2\pi}$

\vspace{1mm}
Is $|a| > 1$, dan heeft $\sin x = a$ geen oplossingen.
\end{stelling}
\begin{bewijs}
Stel $\theta = \bgsin a$, dan is $\sin\theta = a$. Ook $\sin(\pi - \theta) = \sin\theta = a$ (supplementaire hoeken), en een veelvoud van $2\pi$ erbij verandert de sinus niet. Al die $x$ zijn dus oplossingen.

Omgekeerd: is $\sin x = a$, dan ligt het punt $P$ bij $x$ op de cirkel en op de rechte $y = a$. Die snijpunten zijn $\bigl(\pm\sqrt{1 - a^2}, a\bigr)$. Omdat $-\frac{\pi}{2} \leqslant \theta \leqslant \frac{\pi}{2}$, is $\cos\theta \geqslant 0$: het punt bij $\theta$ is het rechtse snijpunt, en dat bij $\pi - \theta$ het linkse. Dus hoort $x$ bij hetzelfde punt als $\theta$ of als $\pi - \theta$, en verschilt het daarvan een veelvoud van $2\pi$.

Is $|a| > 1$, dan snijdt de rechte $y = a$ de cirkel niet: de sinus ligt altijd tussen $-1$ en $1$.
\end{bewijs}
\end{theorie}

Voor $a = 1$ en $a = -1$ raakt de rechte de cirkel, en vallen de twee punten samen. Dan is er maar één reeks oplossingen, zoals $\sin x = 1 \iff x = \frac{\pi}{2} + k \cdot 2\pi$.

\voorbeeld{}{Los op: $\sin x = \frac12$. Geef de algemene oplossing en de oplossingen in $[0, 2\pi[$.\label{vb:h4-sin}}{%
\uitleg{$\bgsin\frac12 = \frac{\pi}{6}$ is een bijzondere waarde; de tweede oplossing is de supplementaire hoek.}
\begin{stappenbreed}
x = \tfrac{\pi}{6} + k \cdot 2\pi \ \text{ of } \ x = \tfrac{5\pi}{6} + k \cdot 2\pi & $\pi - \frac{\pi}{6} = \frac{5\pi}{6}$, $k \in \Z$. \\[1ex]
x \in \bigl\{\tfrac{\pi}{6}, \tfrac{5\pi}{6}\bigr\} & In $[0, 2\pi[$: $k = 0$. Voor $k = 1$ of $k = -1$ val je buiten het interval.
\end{stappenbreed}}

\needspace{6\baselineskip}
\begin{oefeningen}
\oef{h4-sin}
\end{oefeningen}

\needspace{14\baselineskip}
\subsection{De vergelijking $\cos x = a$}

\naast{%
Voor de cosinus snijd je de cirkel met een verticale rechte $x = a$. De twee snijpunten zijn elkaars spiegelbeeld om de $x$-as. Het bovenste punt hoort bij $\theta = \bgcos a$, het onderste bij de tegengestelde hoek $-\theta$.}{%
\begin{tikzpicture}[scale=1.3,font=\scriptsize]
  \eenheidscirkel
  \draw[accent] (.4,-1.2) -- (.4,1.25) node[above] {$x = a$};
  \draw[accent] (0,0) -- (66.4:1) (0,0) -- (-66.4:1);
  \fill[accent] (66.4:1) circle (.035);
  \fill[accent] (-66.4:1) circle (.035);
  \node[accent] at (30:.45) {$\theta$};
  \node[accent] at (-30:.48) {$-\theta$};
\end{tikzpicture}}

\begin{theorie}
\begin{stelling}{De vergelijking $\cos x = a$}
Voor $-1 \leqslant a \leqslant 1$ en met $k \in \Z$:

\vspace{1mm}
\qquad $\formule{\cos x = a \iff x = \bgcos a + k \cdot 2\pi \ \text{ of } \ x = -\bgcos a + k \cdot 2\pi}$

\vspace{1mm}
Is $|a| > 1$, dan heeft $\cos x = a$ geen oplossingen.
\end{stelling}
\begin{bewijs}
Stel $\theta = \bgcos a$. Dan is $\cos\theta = a$, en ook $\cos(-\theta) = \cos\theta = a$ (tegengestelde hoeken). Veelvouden van $2\pi$ erbij veranderen de cosinus niet.

Omgekeerd: is $\cos x = a$, dan ligt het punt bij $x$ op de rechte $x = a$, dus in $\bigl(a, \pm\sqrt{1 - a^2}\bigr)$. Omdat $0 \leqslant \theta \leqslant \pi$, is $\sin\theta \geqslant 0$: het punt bij $\theta$ is het bovenste snijpunt, en dat bij $-\theta$ het onderste. Dus verschilt $x$ een veelvoud van $2\pi$ van $\theta$ of van $-\theta$. Voor $|a| > 1$ snijdt de rechte de cirkel niet.
\end{bewijs}
\end{theorie}

\voorbeeld{}{Los op in $[0, 2\pi[$, op vier decimalen: $\cos x = 0{,}6$.\label{vb:h4-cosel}}{%
\naast{%
\uitleg{Het rekentoestel geeft $\bgcos 0{,}6$. De tweede reeks is de tegengestelde hoek. Om in $[0, 2\pi[$ te belanden, tellen we bij $-0{,}9273$ één keer $2\pi$ op.}
\begin{stappenvrij}{4.4cm}{3.4cm}
\bgcos 0{,}6 \approx 0{,}9273 & Rekentoestel in radialen. \\[1ex]
x \approx \pm 0{,}9273 + k \cdot 2\pi & Algemene oplossing. \\[1ex]
x \approx 0{,}9273 & $+$, met $k = 0$. \\[1ex]
x \approx -0{,}9273 + 2\pi \approx 5{,}3559 & $-$, met $k = 1$.
\end{stappenvrij}}{%
\begin{tikzpicture}[scale=1.25,font=\scriptsize]
  \eenheidscirkel
  \draw[accent] (.6,-1.15) -- (.6,1.2);
  \draw[accent] (0,0) -- (53.13:1) (0,0) -- (-53.13:1);
  \fill[accent] (53.13:1) circle (.035) node[above right] {$0{,}9273$};
  \fill[accent] (-53.13:1) circle (.035) node[below right] {$5{,}3559$};
  \node[below left,inner sep=1pt] at (.6,0) {$0{,}6$};
\end{tikzpicture}}}

\needspace{6\baselineskip}
\begin{oefeningen}
\oef{h4-cos-k}
\end{oefeningen}

\needspace{14\baselineskip}
\subsection{De vergelijking $\tan x = a$}

De tangens heeft periode $\pi$, en neemt op elke tak elke waarde één keer aan. Daarom is er maar één reeks oplossingen, met stappen van $\pi$ in plaats van $2\pi$.

\begin{theorie}
\begin{stelling}{De vergelijking $\tan x = a$}
Voor elke $a \in \R$ en met $k \in \Z$: \quad $\formule{\tan x = a \iff x = \bgtan a + k\pi}$
\end{stelling}
\begin{bewijs}
Stel $\theta = \bgtan a$. Dan is $\tan\theta = a$, en omdat de tangens periode $\pi$ heeft, ook $\tan(\theta + k\pi) = a$.

Omgekeerd: is $\tan x = a$, dan ligt $x$ niet op een asymptoot, en is er een $k \in \Z$ met $x - k\pi \in \bigl]-\frac{\pi}{2}, \frac{\pi}{2}\bigr[$. Ook $\tan(x - k\pi) = a$. De punten bij $x - k\pi$ en bij $\theta$ liggen allebei rechts van de $y$-as, op een rechte door $O$ met richtingscoëfficiënt $a$. Dat is één en dezelfde rechte, die de cirkel rechts maar in één punt snijdt. Twee hoeken in $\bigl]-\frac{\pi}{2}, \frac{\pi}{2}\bigr[$ bij hetzelfde punt zijn gelijk, dus $x - k\pi = \theta$.
\end{bewijs}
\end{theorie}

\voorbeeld{}{Los op in $[0, 2\pi[$, op vier decimalen: $\tan x = 2$.\label{vb:h4-tanel}}{%
\uitleg{Eén reeks oplossingen, met stappen van $\pi$. In $[0, 2\pi[$ passen twee stappen.}
\begin{stappenbreed}
x \approx 1{,}1071 + k\pi & $\bgtan 2 \approx 1{,}1071$. \\[1ex]
x \approx 1{,}1071 \ \text{ of } \ x \approx 1{,}1071 + \pi \approx 4{,}2487 & $k = 0$ en $k = 1$.
\end{stappenbreed}
Controle op de eenheidscirkel: de tweede hoek is antisupplementair met de eerste, en die hebben dezelfde tangens.}

\needspace{6\baselineskip}
\begin{oefeningen}
\oef{h4-tan-k}
\end{oefeningen}

\needspace{12\baselineskip}
\subsection{Een veelvoud van $x$}

Staat er een veelvoud van $x$ in de sinus of cosinus, los dan eerst op naar dat veelvoud, en deel daarna, \textit{ook de periode}. Doe je dat niet, dan mis je oplossingen.

\voorbeeld{}{Los op: $\cos 2x = -\frac12$. Geef de oplossingen in $[0, 2\pi[$.\label{vb:h4-cos}}{%
\uitleg{Los op naar $2x$, en deel dan door 2. De periode $2\pi$ wordt $\pi$: er zijn dus twee keer zoveel oplossingen in $[0, 2\pi[$.}
\begin{stappenbreed}
2x = \tfrac{2\pi}{3} + k \cdot 2\pi \ \text{ of } \ 2x = -\tfrac{2\pi}{3} + k \cdot 2\pi & $\bgcos\bigl(-\frac12\bigr) = \frac{2\pi}{3}$. \\[1ex]
x = \tfrac{\pi}{3} + k\pi \ \text{ of } \ x = -\tfrac{\pi}{3} + k\pi & Delen door 2, ook $k \cdot 2\pi$. \\[1ex]
x \in \bigl\{\tfrac{\pi}{3}, \tfrac{2\pi}{3}, \tfrac{4\pi}{3}, \tfrac{5\pi}{3}\bigr\} & $k = 0, 1$ en $k = 1, 2$.
\end{stappenbreed}
Wie eerst door 2 deelt en pas daarna $k \cdot 2\pi$ optelt, vindt alleen $\frac{\pi}{3}$ en $\frac{5\pi}{3}$: twee van de vier oplossingen.}

\begin{werkwijze}
Zonder $\sin$, $\cos$ of $\tan$ af. Zoek één oplossing met de tabel of met $\bgsin$, $\bgcos$, $\bgtan$ (rekentoestel in radialen). Schrijf de tweede oplossing met de verwante hoek en tel $k \cdot 2\pi$ op (bij de tangens: $k\pi$). Los daarna op naar $x$. Zoek de oplossingen in het gevraagde interval door $k = 0, 1, -1, \ldots$ in te vullen.
\end{werkwijze}

\begin{oefeningen}
\oef{h4-elementair}

\needspace{5\baselineskip}
\oef{h4-veelvoud}
\end{oefeningen}

%========================================================================================
\section{Herleidbare vergelijkingen}

Andere vergelijkingen breng je met de formules uit hoofdstuk 1 en 2 terug tot elementaire vergelijkingen. Er zijn drie veelgebruikte manieren: substitutie, ontbinden en samenvoegen.

\subsection{Substitutie}

Staat er maar één goniometrische functie in de vergelijking, zoals $\sin x$, dan kun je die als nieuwe onbekende nemen. Je krijgt een veeltermvergelijking die je al kunt oplossen.

\voorbeeld{}{Los op: $2\sin^2 x - \sin x - 1 = 0$.\label{vb:h4-kwadratisch}}{%
\uitleg{Een tweedegraadsvergelijking in $\sin x$. Stel $t = \sin x$.}
\begin{stappenbreed}
2t^2 - t - 1 = 0 \iff t = 1 \text{ of } t = -\tfrac12 & Discriminant $9$. Beide liggen in $[-1, 1]$. \\[1ex]
\sin x = 1 \iff x = \tfrac{\pi}{2} + k \cdot 2\pi & De rechte $y = 1$ raakt de cirkel: één reeks. \\[1ex]
\sin x = -\tfrac12 \iff x = -\tfrac{\pi}{6} + k \cdot 2\pi \text{ of } x = \tfrac{7\pi}{6} + k \cdot 2\pi & $\pi - \bigl(-\frac{\pi}{6}\bigr) = \frac{7\pi}{6}$.
\end{stappenbreed}}

Staan er twee functies, zoals $\sin^2 x$ en $\cos x$, schrijf dan eerst alles met één functie. De grondformule of een verdubbelingsformule helpt daarbij.

\voorbeeld{}{Los op: $2\sin^2 x + 3\cos x - 3 = 0$.\label{vb:h4-grondformule}}{%
\uitleg{Er staan een sinus en een cosinus. Omdat de sinus in het kwadraat staat, vervangen we $\sin^2 x$ door $1 - \cos^2 x$. Dan staat er alleen nog $\cos x$.}
\begin{stappenbreed}
2(1 - \cos^2 x) + 3\cos x - 3 = 0 & Grondformule. \\[1ex]
2\cos^2 x - 3\cos x + 1 = 0 & Uitwerken en vermenigvuldigen met $-1$. \\[1ex]
2t^2 - 3t + 1 = 0 \iff t = 1 \text{ of } t = \tfrac12 & Stel $t = \cos x$. Discriminant $1$. \\[1ex]
\cos x = 1 \iff x = k \cdot 2\pi & Eén reeks: $\bgcos 1 = 0$ en $-0 = 0$. \\[1ex]
\cos x = \tfrac12 \iff x = \pm\tfrac{\pi}{3} + k \cdot 2\pi &
\end{stappenbreed}}

Vind je bij de substitutie een $t$ buiten $[-1, 1]$, dan geeft die geen oplossingen: een sinus of cosinus kan zo'n waarde niet aannemen.

\begin{werkwijze}
Schrijf alles met één goniometrische functie, met de grondformule of een verdubbelingsformule. Stel $t$ gelijk aan die functie, en los de veeltermvergelijking op. Schrap de waarden van $t$ buiten $[-1, 1]$ (niet bij de tangens). Los voor elke overblijvende $t$ de elementaire vergelijking op.
\end{werkwijze}

\needspace{5\baselineskip}
\begin{oefeningen}
\oef{h4-kwadr}
\end{oefeningen}

\needspace{14\baselineskip}
\subsection{Ontbinden}

Staat er in elke term eenzelfde factor, breng dan alles naar één lid en zet die factor buiten haakjes. Een product is nul als een van de factoren nul is.

\voorbeeld{}{Los op: $\sin 2x = \cos x$.\label{vb:h4-ontbinden}}{%
\uitleg{Schrijf $\sin 2x$ met de verdubbelingsformule. Dan staat $\cos x$ in beide leden: breng alles naar één lid en ontbind. Deel niet door $\cos x$: dan verlies je oplossingen.}
\begin{stappenbreed}
2\sin x\cos x - \cos x = 0 \iff \cos x\,(2\sin x - 1) = 0 & Gemeenschappelijke factor. \\[1ex]
\cos x = 0 \iff x = \tfrac{\pi}{2} + k\pi & Eerste factor nul. \\[1ex]
\sin x = \tfrac12 \iff x = \tfrac{\pi}{6} + k \cdot 2\pi \text{ of } x = \tfrac{5\pi}{6} + k \cdot 2\pi & Tweede factor nul.
\end{stappenbreed}}

Waarom niet delen door $\cos x$? Wie deelt, krijgt $2\sin x = 1$, en vindt alleen de tweede reeks. Maar ook $x = \frac{\pi}{2}$ is een oplossing: $\sin\pi = 0 = \cos\frac{\pi}{2}$. Delen door een uitdrukking mag alleen als je zeker weet dat ze niet nul is.

\begin{werkwijze}
Breng alles naar één lid. Zet een gemeenschappelijke factor buiten haakjes, eventueel na een verdubbelingsformule. Stel elke factor gelijk aan nul, en los elke elementaire vergelijking op.
\end{werkwijze}

\needspace{5\baselineskip}
\begin{oefeningen}
\oef{h4-ontbind}
\end{oefeningen}

\needspace{14\baselineskip}
\subsection{Samenvoegen}

Staat er een som $a\sin x + b\cos x$ gelijk aan een getal, dan helpt substitutie niet: er staan twee functies, zonder kwadraat. Voeg ze samen tot één sinus, zoals in hoofdstuk 2.

\voorbeeld{}{Los op: $\sin x + \sqrt3\cos x = 1$.\label{vb:h4-samenvoegen}}{%
\uitleg{Voeg samen tot één sinus: $\sin x + \sqrt3\cos x = 2\sin\bigl(x + \frac{\pi}{3}\bigr)$. Los daarna op naar $x + \frac{\pi}{3}$.}
\begin{stappenbreed}
\sin\bigl(x + \tfrac{\pi}{3}\bigr) = \tfrac12 & Delen door $2$. \\[1ex]
x + \tfrac{\pi}{3} = \tfrac{\pi}{6} + k \cdot 2\pi \ \text{ of } \ x + \tfrac{\pi}{3} = \tfrac{5\pi}{6} + k \cdot 2\pi & Elementaire vergelijking. \\[1ex]
x = -\tfrac{\pi}{6} + k \cdot 2\pi \ \text{ of } \ x = \tfrac{\pi}{2} + k \cdot 2\pi & Controle: $\sin\frac{\pi}{2} + \sqrt3\cos\frac{\pi}{2} = 1$.
\end{stappenbreed}}

\begin{oefeningen}
\oef{h4-samen}
\end{oefeningen}

\subsection{Welke methode?}

Kijk eerst welke functies er staan. Eén functie, eventueel in het kwadraat: substitutie. Een gemeenschappelijke factor: ontbinden. Een som van een sinus en een cosinus van dezelfde hoek, zonder kwadraten: samenvoegen.

\needspace{5\baselineskip}
\begin{oefeningen}
\oef{h4-herleidbaar}
\end{oefeningen}

%========================================================================================
\section{Vergelijkingen in een context}

In een context vraag je vaak niet wanneer een grootheid een waarde \textit{is}, maar wanneer ze boven of onder die waarde ligt. Dan los je eerst de vergelijking op, en lees je daarna op de grafiek af waar de functie boven of onder de rechte ligt.

\voorbeeld{}{In de haven uit hoofdstuk 3 is de waterhoogte $h(t) = 2{,}2\cos\bigl(\frac{2\pi}{12{,}5}(t - 2)\bigr) + 2{,}6$ (in m, $t$ in uur na middernacht). Een schip heeft minstens $3$ m water nodig. Wanneer kan het tussen middernacht en 18 uur binnenvaren?\label{vb:h4-context}}{%
\uitleg{Los eerst de vergelijking $h(t) = 3$ op. Rond de top (om 2 uur) is het water hoger dan $3$ m; de grenzen liggen symmetrisch rond die top.}
\begin{stappenbreed}
\cos\bigl(\tfrac{2\pi}{12{,}5}(t - 2)\bigr) = \tfrac{0{,}4}{2{,}2} \approx 0{,}1818 & Min $2{,}6$, en delen door $2{,}2$. \\[1ex]
\tfrac{2\pi}{12{,}5}(t - 2) = \pm 1{,}388 & $\bgcos 0{,}1818 \approx 1{,}388$, met $k = 0$. \\[1ex]
t - 2 = \pm 2{,}76 & Vermenigvuldigen met $\frac{12{,}5}{2\pi}$. \\[1ex]
t \in [0;\ 4{,}76] \text{ en } [11{,}74;\ 17{,}26] & Tel er een periode $12{,}5$ bij voor de tweede.
\end{stappenbreed}
\figuren{\begin{tikzpicture}[x=.6cm,y=.45cm,font=\scriptsize]
  \fill[accent!15] (0,0) rectangle (4.76,5);
  \fill[accent!15] (11.74,0) rectangle (17.26,5);
  \draw[black!14,very thin] (0,0) grid[xstep=2,ystep=1] (18,5);
  \draw[->,grijs] (0,0) -- (18.6,0) node[right,text=black] {$t$};
  \draw[->,grijs] (0,0) -- (0,5.6) node[above,text=black] {$h$};
  \foreach \t in {2,6,10,14,18}{\node[below,inner sep=1.5pt] at (\t,0) {$\t$};}
  \foreach \t in {1,3,5}{\node[left,inner sep=1.5pt] at (0,\t) {$\t$};}
  \draw[grijs,thick] (0,3) -- (18,3);
  \draw[accent,thick] plot[domain=0:18,samples=120,smooth] (\x,{2.2*cos(2*pi/12.5*(\x-2) r)+2.6});
  \foreach \p in {4.76,11.74,17.26}{\fill[accent] (\p,3) circle (2.2pt);}
\end{tikzpicture}}
Het schip kan binnenvaren tot ongeveer 4.45 uur, en opnieuw tussen ongeveer 11.45 en 17.15 uur.}

\begin{werkwijze}
Vervang het ongelijkheidsteken door een gelijkheidsteken en los die vergelijking op in het gevraagde interval. Schets de grafiek met de horizontale rechte, en lees af tussen welke oplossingen de functie boven (of onder) de rechte ligt. Geef het antwoord in de eenheden van de context.
\end{werkwijze}

\begin{oefeningen}
\oef{h4-context}

\needspace{7\baselineskip}
\oef{h4-schommel}
\end{oefeningen}

%========================================================================================
\opdrachtenkop

\begin{oefeningen}
\needspace{10\baselineskip}
\oef{h4-o1}

\needspace{8\baselineskip}
\oef{h4-o10}

\needspace{14\baselineskip}
\oef{h4-o12}

\needspace{8\baselineskip}
\oef{h4-o2}

\needspace{12\baselineskip}
\oef{h4-o13}

\needspace{8\baselineskip}
\oef{h4-o15}

\needspace{6\baselineskip}
\oef{h4-o3}

\needspace{9\baselineskip}
\oef{h4-o11}

\needspace{10\baselineskip}
\oef{h4-o14}

\needspace{12\baselineskip}
\oef{h4-o16}

\needspace{12\baselineskip}
\oef{h4-o20}

\oef{h4-o4}

\needspace{14\baselineskip}
\oef{h4-o22}

\needspace{8\baselineskip}
\oef{h4-o18}

\needspace{14\baselineskip}
\oef{h4-o19}

\needspace{9\baselineskip}
\oef{h4-o17}

\oef{h4-o7}

\oef{h4-o5}

\needspace{14\baselineskip}
\oef{h4-o23}

\oef{h4-o6}

\needspace{14\baselineskip}
\oef{h4-o21}

\needspace{12\baselineskip}
\oef{h4-o24}

\needspace{10\baselineskip}
\oef{h4-o9}

\needspace{12\baselineskip}
\oef{h4-o25}

\oef{h4-o8}

\needspace{16\baselineskip}
\oef{h4-o27}

\needspace{12\baselineskip}
\oef{h4-o26}
\end{oefeningen}

%----------------------------------------------------------------------------------------
% KOPPEL: verbanden met andere leerstof
%----------------------------------------------------------------------------------------
\needspace{20\baselineskip}
\opdrachtsectie{Koppel}

\lkr{Algemene functies: kijk na wat al gezien is. Het pakket legt de volgorde niet vast.}

\begin{oefeningen}

\oef{h4-koppel-2}

\end{oefeningen}

%----------------------------------------------------------------------------------------
\opdrachtsectie{Test jezelf}

\begin{oefeningen}
\oef{h4-t1}

\oef{h4-t2}

\oef{h4-t3}

\oef{h4-t4}

\oef{h4-t6}

\oef{h4-t7}

\oef{h4-t5}
\end{oefeningen}

\herhaalschema{%
  \fv{oef:h4-t1}{vb:h4-bg}{oef:h4-bg}%
  \fv{oef:h4-t2}{vb:h4-cosel}{oef:h4-cos-k}%
  \fv{oef:h4-t3}{vb:h4-cos}{oef:h4-veelvoud}%
  \fv{oef:h4-t4}{vb:h4-ontbinden}{oef:h4-ontbind}%
  \fv{oef:h4-t6}{vb:h4-grondformule}{oef:h4-kwadr}%
  \fv{oef:h4-t7}{vb:h4-samenvoegen}{oef:h4-samen}%
  \fv{oef:h4-t5}{vb:h4-context}{oef:h4-context}}

\ifzelfstudie
\section*{Maak je eigen samenvatting}
Sluit dit boek en schrijf in je schrift: domein en bereik van $\bgsin$, $\bgcos$ en $\bgtan$, de algemene oplossing van $\sin x = a$, $\cos x = a$ en $\tan x = a$ met een figuur van de eenheidscirkel, hoe je oplossingen in een interval vindt, en de drie manieren om een vergelijking te herleiden (substitutie, ontbinden, samenvoegen). Vergelijk daarna met de kaders.
\fi
